The comparisons \pyinline{==}, \pyinline{!=}, \pyinline{in} and \pyinline{not in} compare values
structurally\speconly{, by structural equality $\eqOp{\Sigma}$ and membership $\inOp{\Sigma}$}, and agree with Python
except at the comparisons left undefined by \PurePy\paperonly{, described in \secref{no-coercions}}\speconly{
(\secref{no-coercions})}: between a Boolean and a number, between two functions, and between two \pyinline{nan}
values inside a container. A run that reaches one of these comparisons has undefined behaviour, so an
implementation may abort or produce the result given by Python (\paperonly{\secref{python-compatibility}}\speconly{\secref{conformance}}).
Equality and membership compare the elements of a sequence in order and stop at the first element that decides
the result, so an undefined comparison later in the sequence does not make the whole comparison undefined. Membership in a range, or of a key other than a string in a dictionary, is ruled out by typing and aborts with $\errExcluded$\speconly{ (\S\ref{aux:in})}.
